\section{Introduction}

One of the goals of lattice calculations is to predict the low-energy hadron
spectrum of confined quarks and gluons, starting solely from the QCD
lagrangian. This approach to spectroscopy necessitates methods for measuring the
two-point correlation functions of field operators with the selected quantum
numbers under investigation. A complete understanding of the QCD spectrum
must also included excitations of the mesons and baryons as well as exotic
states.
Describing the resonances seen in scattering experiments
as states in QCD has long presented a challenge to lattice calculations as
direct access to the matrix elements related to decay widths is usually
missing in Euclidean formulations of field theory. The way to circumvent this
is well known in principle: decay properties can be inferred from a detailed
study of the dependence of the spectrum in a finite box on the volume
of the box \cite{DeWitt:1956be,Luscher:1991cf}.
The best means of making solid determinations of energy levels is to employ a
large basis of operators and then use the variational method to find good
approximations to energy eigenstates. As soon as states above threshold are to
be explored, the basis should include operators that resemble multi-hadron
systems where the constituents have well-defined momenta.
Access to all elements of the quark propagator \cite{Foley:2005ac} relevant
to long-range correlation enables these measurements and
also enables more detailed investigations into physical observables
whose accurate measurement has often eluded lattice practitioners. One example
is the isoscalar meson sector, where mass determinations require the
evaluation of disconnected diagrams as part of the two-point correlation
function.  As well as
this ambitious list of requirements, precision data
will be crucial for these calculations \cite{Morningstar:2008mc}.

In a Monte Carlo calculation on the lattice, the physically relevant signal in
a correlation function falls exponentially and is rapidly dominated by
statistical fluctuations. Operators that create low-lying energy eigenstates
quickly are invaluable and improve the quality of data extracted
exponentially. The most useful tool at the lattice practitioner's disposal in
building good creation operators is smearing. The best means of
understanding the low-energy degrees of freedom of confinement must thus come
from a smearing method that can address all these features; it must be
numerically accessible, do a good job of projecting onto the lowest energy
states and it must facilitate easy evaluation of correlation functions
involving a broad set of creation operators, including those exciting
multi-hadron states.

A new generation of experiments devoted to hadron spectroscopy,
including GlueX at Jefferson Lab, PANDA at GSI/FAIR and BES III
intend to make measurements with unprecedented precision and in
previously unexplored mass ranges and quantum-number sectors. The main aim of
the {\em Hadron Spectrum Collaboration} is to extract from realistic
lattice simulations predictions for masses, decay properties and relevant matrix
elements in these domains.

Spectroscopy investigations are almost as old as non-perturbative Monte Carlo
calculations on the lattice \cite{Hamber:1981zn,Weingarten:1980hx}.  As the
first studies progressed, it was quickly realised that operators
creating hadrons which are constructed directly from the fields in the lattice
lagrangian have significant overlap with a large tower of states and that
extracting the lightest elements of the spectrum was difficult. The solution was
to build operators from ``smeared'' fields, where some linear operator was
applied first to the quark degrees of freedom on the appropriate time-slice
before the creation operator was formed. The aim of constructing a smeared
field was to reduce the component of the fields close to the
cut-off substantially, since these modes do not contribute significantly to
long-range correlation functions.
Initially, the techniques considered smeared fields that were not
gauge covariant. Gauge-covariant schemes were introduced soon after
\cite{Gusken:1989ad,
Allton:1993wc}. In particular, Ref.~\cite{Allton:1993wc} introduced the Jacobi
smearing algorithm, in which a lattice approximation to the three-dimensional
gauge-covariant laplacian is constructed and applied iteratively to the field.
This acts naturally as a long-wavelength filter on the modes constituting the
quark field.


In this paper, a new method for smearing quark fields is described and tested in
large-scale realistic simulations.
In Sec.~\ref{sec:theory}, the formulation of the new algorithm is presented,
including detail on its application to measurements of hadron two-point
correlation functions and three-point functions where an arbitrary current is
inserted.
Sec.~\ref{sec:results} presents the results of initial tests of the
effectiveness of the method, and Sec.~\ref{sec:discussion} discusses practical
issues that arise with the method and outlines some future research directions.
